% Continuation of the codes section (\cref{sec:codes}): subsections only.

\subsection{Bounded complex coefficients: Theorem A}\label{sec:ana-A}

The rate bounds proved so far in this section hold either for all codes on special graphs (\cref{thm:H} and its
corollaries) or for all linear codes on triangle-free graphs of maximum degree~$3$ (\cref{thm:S}). The undirected crux
(\cref{prob:undirected-crux}) concerns locally bipartite side information of bounded degree. In this subsection we
prove the bound~$\frac12+o(1)$ on the rate in that setting for one class of codes: linear codes over~$\R$ or~$\C$
whose recovery coefficients are bounded.

Let $G$ be a finite graph with vertex set $V$ and $N=\card V$. A matrix $M\in F^{V\times V}$ over a field $F$
\emph{fits} $G$ if $M_{vv}=1$ for every $v$ and $M_{vu}=0$ whenever $u\ne v$ and $uv\notin E(G)$. By
\cref{lem:linear-pattern}, the linear storage codes on $G$ over $F$ are the subspaces of kernels of matrices that fit
$G$; the coefficients $R_{vu}$ of a recovery rule $x_v=\sum_{u\sim v}R_{vu}x_u$ are the negated off-diagonal entries of
such a matrix, which we call a \emph{check matrix} of the code.
A complex matrix is \emph{$B$-bounded} if all its entries have modulus at most $B$; a matrix that fits a
graph can be $B$-bounded only if $B\ge1$. For $\ell\ge1$ let $\mathrm{odd}_\ell(G)$ be the set of vertices of $G$ that
lie on an odd cycle of length at most $\ell$. As in \cref{prob:undirected-crux}, a sequence of graphs $G_n$ on
$N_n\to\infty$ vertices is \emph{locally bipartite} if $\card{\mathrm{odd}_\ell(G_n)}=o(N_n)$ for every fixed $\ell$.
Sequences whose girth tends to infinity
are locally bipartite, and so are \emph{locally tree-like} sequences, in which, for every fixed $\ell$, only $o(N_n)$
vertices lie on cycles of length at most $\ell$.

\begin{theorem}[Theorem A: bounded complex coefficients]\label{thm:A}
\notformalized
Let $D\ge1$ and $B\ge1$.
\begin{enumerate}[label=(\roman*)]
\item Let $(G_n)$ be a locally bipartite sequence of graphs of maximum degree at most $D$ on $N_n$ vertices, and let
  $M_n\in\C^{V(G_n)\times V(G_n)}$ be $B$-bounded matrices that fit $G_n$. Then
  $\liminf_n\rank(M_n)/N_n\ge\frac12$.
\item \emph{(Uniform form.)} For every $\eps>0$ there are $\delta>0$ and $N_0$ such that $\rank M\ge(\frac12-\eps)N$
  for every graph $G$ on $N\ge N_0$ vertices of maximum degree at most $D$ with
  $\card{\mathrm{odd}_{\lfloor1/\delta\rfloor}(G)}\le\delta N$ and every $B$-bounded complex matrix $M$ that fits $G$.
\item \emph{(Random regular graphs.)} Let $d\ge3$ and $\eps>0$, and let $G_{n,d}$ be a uniformly random $d$-regular
  graph on $n$ vertices ($nd$ even). With probability tending to $1$ as $n\to\infty$, every $B$-bounded complex matrix
  that fits $G_{n,d}$ has rank at least $(\frac12-\eps)n$.
\end{enumerate}
\end{theorem}

In terms of codes (\cref{cor:ana-codes}): on a locally bipartite sequence of graphs of bounded degree, a linear storage
code over $\R$ or $\C$ whose recovery coefficients have modulus at most $B$ has dimension at most $(\frac12+o(1))N_n$.

\paragraph{Why it matters.} For the impossibility side, Theorem~A proves the analogue of the crux's rate bound
(\cref{def:crux}) for real and complex linear codes with bounded recovery coefficients, scalar or with a fixed number of
symbols per vertex (\cref{prop:ana-A-ext}), on side information of bounded total degree whose underlying graphs are
locally bipartite. The crux itself concerns binary codes; Theorem~A bears on them through \cref{cor:torsion}. Only the support of the check matrix enters, so directed side information is
covered through its underlying graph, and locally tree-like side information is covered. For the construction side it
is a design rule: on such side information, a code that beats rate~$\frac12$ by a constant is nonlinear, or vector-linear
with a number of symbols per vertex that grows with $N$, or linear with coefficients that no diagonal rescaling makes
bounded (\cref{prop:ana-A-ext}(a), for $t=1$), or linear over a finite field, and then every bounded integer lift of its
check matrix has $p$-torsion of linear rank (\cref{cor:torsion}). Codes over finite alphabets, linear or not, are not
covered; they are the open part of the crux.

\paragraph{Local limits.} We use the framework of Benjamini and Schramm~\cite{BenjaminiSchramm2001} and of Aldous and
Lyons~\cite{AldousLyons2007}. A \emph{network} is a connected graph of maximum degree at most $D$ with \emph{marks}: a
point of a compact metric space $\Xi_V$ at every vertex and a point of a compact metric space $\Xi_E$ at each end of
every edge. A \emph{rooted} network has a distinguished vertex $o$, a \emph{doubly rooted} network an ordered pair of
distinguished vertices, and we identify networks that are isomorphic by graph isomorphisms preserving the roots and the
marks. Two rooted networks are close if, for a large $R$, the balls of radius $R$ around their roots are isomorphic by a
root-preserving graph isomorphism that changes every mark by less than $1/R$. Since degrees are bounded and the mark
spaces are compact, the space of rooted networks is compact (a diagonal argument). A finite marked graph of maximum
degree at most $D$, connected or not, defines a random rooted network: choose $o$ uniformly among its vertices and keep
the component of $o$. A sequence of finite marked graphs \emph{converges locally} to a probability measure $\mu$ on
rooted networks if these random rooted networks converge to $\mu$ in distribution. A probability measure $\mu$ on
rooted networks is \emph{unimodular} if it satisfies the mass-transport principle
\begin{equation}\label{eq:ana-mtp}
  \mathbb{E}_\mu\sum_{x\in V(G)}f(G,o,x)\;=\;\mathbb{E}_\mu\sum_{x\in V(G)}f(G,x,o)
\end{equation}
for every Borel function $f\ge0$ on doubly rooted networks, where $(G,o)$ denotes a $\mu$-random rooted network. We use
four standard facts.
\begin{enumerate}[label=(L\arabic*)]
\item Every sequence of finite marked graphs as above has a locally convergent subsequence (compactness and
  Prokhorov's theorem).
\item Every local limit of finite marked graphs is unimodular~\cite[Section~2]{AldousLyons2007}. (A finite graph with
  several components is a convex combination of its components, each rooted uniformly.)
\item Let $\mu$ be unimodular. A family $T=(T_G)$ of bounded operators $T_G$ on $\ell^2(V(G))$, indexed by networks,
  is \emph{equivariant} if $\langle T_{G'}\delta_{\varphi x},\delta_{\varphi y}\rangle=\langle T_G\delta_x,\delta_y\rangle$
  for every isomorphism $\varphi\colon G\to G'$ of networks and all vertices $x,y$, and \emph{measurable} if
  $(G,x,y)\mapsto\langle T_G\delta_x,\delta_y\rangle$ is Borel on doubly rooted networks. The $\mu$-classes of
  equivariant, measurable families with $\sup_G\lVert T_G\rVert<\infty$ form a von Neumann algebra $\mathrm{Alg}(\mu)$,
  acting on the direct integral of the spaces $\ell^2(V(G))$, on which
  $\mathrm{Tr}(T)=\mathbb{E}_\mu\langle T_G\delta_o,\delta_o\rangle$ is a positive linear functional with
  $\mathrm{Tr}(ST)=\mathrm{Tr}(TS)$~\cite[Section~5]{AldousLyons2007}. Its elements act fibrewise, and their kernels,
  ranges and spectral projections are computed fibrewise.
\item (Polar decomposition.) For $T$ in a von Neumann algebra $\mathcal A$ there is a partial isometry $U\in\mathcal A$
  such that $U^*U$ is the projection onto $(\ker T)^\perp$ and $UU^*$ is the projection onto the closure of the range
  of $T$~\cite[Theorem~6.1.2 and Proposition~6.1.3]{KadisonRingrose1997}.
\end{enumerate}

\begin{informalproof}
(i) Let $\nu_n=1-\rank(M_n)/N_n$ be the normalized nullity and $\nu=\limsup_n\nu_n$; we show $\nu\le\frac12$. Make
$G_n$ a marked graph by putting the mark $M_n[v,u]$, a point of the disc $\{z\in\C:\abs z\le B\}$, at the end $v$ of
each edge $vu$. Pass to a subsequence along which $\nu_n\to\nu$ and, by (L1), the marked graphs converge locally to some
$\mu$, which is unimodular by (L2). Write $(G,M,o)$ for a $\mu$-random rooted network. The marks define an operator $M_G$
on $\ell^2(V(G))$ by $(M_Gx)_v=x_v+\sum_{u\sim v}M[v,u]\,x_u$. Each row and each column of $M_G$ has at most $D+1$
nonzero entries, of modulus at most $B$, so $\lVert M_G\rVert\le(D+1)B$, and $(M_G)$ is an equivariant measurable
family.

\emph{Step 1: the limit is bipartite.} For $\ell,R\ge1$ let $A_{\ell,R}$ be the set of rooted networks in which some odd
cycle of length at most $\ell$ meets the ball of radius $R$ around the root. Such a cycle lies in the ball of radius
$R+\ell$, so membership in $A_{\ell,R}$ is decided by that ball without its marks, and $A_{\ell,R}$ is open and closed.
A uniform root of $G_n$ lies in $A_{\ell,R}$ only if it is within distance $R$ of $\mathrm{odd}_\ell(G_n)$, which has
probability at most $(1+D+\dots+D^R)\card{\mathrm{odd}_\ell(G_n)}/N_n\to0$. Since $A_{\ell,R}$ is open and closed,
$\mu(A_{\ell,R})$ is the limit of these probabilities, so $\mu(A_{\ell,R})=0$. Every odd cycle of $G$ meets some ball
around $o$, so $\mu$-almost surely $G$ has no odd cycle: $G$ is bipartite.

\emph{Step 2: semicontinuity of the nullity.} Let $H_n=M_n^*M_n$ and let $\sigma_n$ be the empirical eigenvalue
distribution of $H_n$ (eigenvalues counted with multiplicity, each with mass $1/N_n$); it is supported in $[0,c]$ with
$c=((D+1)B)^2$. For $k\ge0$, $\int x^k\,d\sigma_n=N_n^{-1}\operatorname{tr}H_n^k$ is the expectation, over a uniform
vertex $o$ of $G_n$, of $\langle H_n^k\delta_o,\delta_o\rangle$. Expanding the product,
$\langle(M^*M)^k\delta_o,\delta_o\rangle$ is a sum of at most $(D+1)^{2k}$ products of $2k$ entries of $M$ or their
conjugates, taken along walks of length $2k$ from $o$ to $o$ whose steps stay put or move to a neighbour; such walks
stay in the ball of radius $k$. Hence $(G,M,o)\mapsto\langle(M_G^*M_G)^k\delta_o,\delta_o\rangle$ is a bounded continuous
function on rooted networks, and $\int x^k\,d\sigma_n\to\mathbb{E}_\mu\langle H^k\delta_o,\delta_o\rangle=\int
x^k\,d\sigma$, where $H=M_G^*M_G$ and $\sigma=\mathbb{E}_\mu\sigma_{H,o}$ is the average of the spectral measures
$\sigma_{H,o}$ of $H$ at $\delta_o$. All these are probability measures on $[0,c]$, so convergence of moments implies
$\sigma_n\to\sigma$ weakly (Weierstrass approximation), and the portmanteau theorem for the closed set $\{0\}$ gives
$\limsup_n\sigma_n(\{0\})\le\sigma(\{0\})$. Since $\ker(M^*M)=\ker M$ for matrices and for bounded operators,
$\sigma_n(\{0\})=\nu_n$ and $\sigma(\{0\})=\mathbb{E}_\mu\langle P_K\delta_o,\delta_o\rangle$, where $P_K$ is the
orthogonal projection onto $K=\ker M_G$. Therefore
\begin{equation}\label{eq:ana-step2}
  \nu\;\le\;\mathbb{E}_\mu\langle P_K\delta_o,\delta_o\rangle .
\end{equation}
This is the easy half of L\"uck's approximation theorem~\cite{Luck1994}; for adjacency matrices of random graphs the
same argument is used by Bordenave, Lelarge and Salez~\cite{BLS2011}. The bound $B$ is used here and in~(L1).

\emph{Step 3: a random colouring.} For $\mu$-almost every $(G,M,o)$ the graph $G$ is connected and bipartite, so it has
exactly two proper colourings with colours $\pm1$, each the negative of the other. Let $\hat\mu$ be the law of
$(G,M,\beta,o)$, where $(G,M,o)$ has law $\mu$ and $\beta$, a vertex mark, is one of the two colourings, chosen by a fair
coin independent of everything else. Then $\hat\mu$ is unimodular. Indeed, let $f\ge0$ be Borel on doubly rooted
coloured networks. For a doubly rooted network $(G,M,x,y)$ with $G$ bipartite let $\pi_x(w)=(-1)^{\dist(x,w)}$, and put
$\bar f(G,M,x,y)=\frac12\bigl(f(G,M,\pi_x,x,y)+f(G,M,-\pi_x,x,y)\bigr)$ (and $\bar f=0$ if $G$ is not bipartite).
Since $\{\pi_x,-\pi_x\}=\{\pi_y,-\pi_y\}$ is the set of both colourings, $\bar f$ is the average of $f$ over the two
colourings whichever root is used to define it; it is Borel and invariant under isomorphisms. Hence
\[
  \mathbb{E}_{\hat\mu}\sum_xf(G,M,\beta,o,x)=\mathbb{E}_\mu\sum_x\bar f(G,M,o,x)
  =\mathbb{E}_\mu\sum_x\bar f(G,M,x,o)=\mathbb{E}_{\hat\mu}\sum_xf(G,M,\beta,x,o),
\]
the middle equality being \eqref{eq:ana-mtp} for $\mu$. Moreover $\hat\mu(\beta(o)=-1)=\frac12$. (Without the coin,
colouring by the parity of the distance to the root is not unimodular in general: the root would always have colour
$+1$ and its neighbours colour $-1$.)

\emph{Step 4: the trace.} By (L3) the algebra $\mathrm{Alg}(\hat\mu)$ contains $M=(M_G)$, hence $H=M^*M$ and its spectral
projection $P_K=\mathbf 1_{\{0\}}(H)$, and the projection $P_J$ onto $\ell^2(J)$, where $J=\beta^{-1}(-1)$; let
$V_+=\beta^{-1}(+1)$. Let $T=P_JP_K\in\mathrm{Alg}(\hat\mu)$. If $Tx=0$, write $x=k+k'$ with $k\in K$ and
$k'\in K^\perp$; then $P_Jk=Tx=0$, so $k$ vanishes on $J$. Every neighbour of a vertex $v\in V_+$ lies in $J$, so
$0=(M_Gk)_v=k_v+\sum_{u\sim v}M[v,u]\,k_u=k_v$. Hence $k=0$, and $\ker T=K^\perp$. The range of $T$ lies in
$\ell^2(J)$. By (L4) there is a partial isometry $U\in\mathrm{Alg}(\hat\mu)$ with $U^*U=P_K$ and $UU^*\le P_J$, and by
(L3)
\[
  \mathbb{E}_\mu\langle P_K\delta_o,\delta_o\rangle=\mathrm{Tr}(P_K)=\mathrm{Tr}(U^*U)=\mathrm{Tr}(UU^*)
  \le\mathrm{Tr}(P_J)=\hat\mu(\beta(o)=-1)=\tfrac12 .
\]
The first equality holds because $\langle P_K\delta_o,\delta_o\rangle$ does not depend on $\beta$ and the law of
$(G,M,o)$ under $\hat\mu$ is $\mu$. With \eqref{eq:ana-step2}, $\nu\le\frac12$.

(ii) If (ii) failed for some $\eps$, then for every $m$ there would be a graph $G_m$ on $N_m\ge m$ vertices of maximum
degree at most $D$ with $\card{\mathrm{odd}_m(G_m)}\le N_m/m$, and a $B$-bounded matrix $M_m$ fitting $G_m$ with
$\rank M_m<(\frac12-\eps)N_m$. For fixed $\ell$ and $m\ge\ell$ we have
$\card{\mathrm{odd}_\ell(G_m)}\le\card{\mathrm{odd}_m(G_m)}\le N_m/m$, so $(G_m)$ is locally bipartite, and (i)
gives a contradiction.

(iii) By a theorem of Wormald~\cite{Wormald1981}, for each fixed $\ell\ge3$ the number $X_\ell$ of cycles of length
$\ell$ in $G_{n,d}$ converges in distribution as $n\to\infty$, to a Poisson law with mean $(d-1)^\ell/(2\ell)$; in
particular $X_\ell$ is bounded in probability. Let $\delta$ and $N_0$ be given by (ii) for $D=d$. Then
$\card{\mathrm{odd}_{\lfloor1/\delta\rfloor}(G_{n,d})}\le\sum_{\ell\le1/\delta}\ell X_\ell$ is bounded in probability,
so it is at most $\delta n$ with probability tending to $1$, and (ii) applies.
\end{informalproof}

\begin{proposition}[Extensions of Theorem A]\label{prop:ana-A-ext}
\notformalized
Let $D\ge1$ and $B\ge1$.
\begin{enumerate}[label=(\alph*)]
\item \emph{(Tight coefficients and gauge.)} For every $\eps>0$ there are $\delta'>0$ and $N_0'$ with the following
  property. Let $G$ be a graph on $N\ge N_0'$ vertices of maximum degree at most $D$ with
  $\card{\mathrm{odd}_{\lfloor1/\delta'\rfloor}(G)}\le\delta'N$, let $M\in\C^{V\times V}$ fit $G$, let $S$ be an
  invertible diagonal matrix, and let $X$ be the set of rows of $S^{-1}MS$ that contain an entry of modulus larger than
  $B$. Then $\rank M\ge(\frac12-\eps)(N-\card X)$.
\item \emph{(A fixed number of symbols per vertex.)} Let $t\ge1$. Parts (i) and (ii) of \cref{thm:A} hold for complex
  matrices $M$ with rows and columns indexed by $V\times[t]$ and entries of modulus at most $B$ whose $t\times t$ blocks
  satisfy $M_{vv}=I_t$ and $M_{vu}=0$ whenever $u\ne v$ and $uv\notin E(G)$, with $N$ replaced by $tN$ in the
  conclusions; the constants $\delta$ and $N_0$ depend on $t$.
\item \emph{(Invariant independent sets.)} Let $B$-bounded matrices $M_n$ fitting graphs $G_n$ of maximum degree at most
  $D$, viewed as marked graphs as in the proof of \cref{thm:A}, converge locally to $\mu$. Let $\alpha_U(\mu)$ be the
  supremum of $\hat\mu(\iota(o)=1)$ over all unimodular laws $\hat\mu$ of networks with an additional vertex mark
  $\iota\in\{0,1\}$ such that forgetting $\iota$ gives $\mu$ and $\iota^{-1}(1)$ is $\hat\mu$-almost surely an
  independent set. Then $\liminf_n\rank(M_n)/N_n\ge\alpha_U(\mu)$.
\item \emph{(Sharpness.)} If $G$ has a perfect matching $P$ with adjacency matrix $A_P$, then $I+A_P$ is $1$-bounded,
  fits $G$, and has rank $N/2$. In particular the constant $\frac12$ of \cref{thm:A} is attained on every regular
  bipartite graph of positive degree.
\end{enumerate}
\end{proposition}

\begin{informalproof}
(a) The matrix $S^{-1}MS$ fits $G$ and has the rank of $M$, so we may assume $S=I$. Let $\delta$ and $N_0$ be given by
\cref{thm:A}(ii) for $D$, $B$ and $\eps$, and put $\delta'=2\delta/(D+1)$ and $N_0'=\lceil(D+1)N_0/2\rceil$. Let
$N'=N-\card X$. If $N'\le2N/(D+1)$, let $Y$ be an independent set of $G$ with $\card Y\ge N/(D+1)$ (a greedy choice);
the principal submatrix $M[Y,Y]$ is the identity, so $\rank M\ge\card Y\ge N'/2$. Otherwise $N'>2N/(D+1)\ge N_0$, and
the principal submatrix $M'=M[V\setminus X,V\setminus X]$ is $B$-bounded and fits $G'=G[V\setminus X]$. Every odd cycle
of $G'$ is an odd cycle of $G$, and $\lfloor1/\delta\rfloor\le\lfloor1/\delta'\rfloor$, so
$\card{\mathrm{odd}_{\lfloor1/\delta\rfloor}(G')}\le\card{\mathrm{odd}_{\lfloor1/\delta'\rfloor}(G)}\le\delta'N<\delta N'$.
\Cref{thm:A}(ii) gives $\rank M\ge\rank M'\ge(\frac12-\eps)N'$.

(b) Let $G[t]$ be the graph on $V\times[t]$ in which $(v,i)\sim(u,j)$ if and only if $uv\in E(G)$. Such a matrix $M$
fits $G[t]$, whose maximum degree is $tD$. An odd cycle of length $\ell'\le\ell$ through $(v,i)$ in $G[t]$ projects to a
closed walk of odd length $\ell'$ through $v$ in $G$; a closed walk of odd length contains an odd cycle of length at
most $\ell'$ formed by vertices of the walk, all within distance $\ell'/2\le\ell$ of $v$. Hence
$\mathrm{odd}_\ell(G[t])\subseteq W\times[t]$, where $W$ is the set of vertices within distance $\ell$ of
$\mathrm{odd}_\ell(G)$, and $\card{\mathrm{odd}_\ell(G[t])}\le t(1+D+\dots+D^\ell)\card{\mathrm{odd}_\ell(G)}$. So
$(G_n[t])$ is locally bipartite when $(G_n)$ is, and \cref{thm:A}(i) applies to $G_n[t]$ with $D$ replaced by $tD$.
For~(ii), take $\delta_t$ and $N_0$ from \cref{thm:A}(ii) for degree $tD$, put $\ell=\lfloor1/\delta_t\rfloor$ and
$\delta=\min\bigl(\delta_t/(1+D+\dots+D^\ell),1/\ell\bigr)$. Then $\card{\mathrm{odd}_{\lfloor1/\delta\rfloor}(G)}\le\delta N$
gives $\card{\mathrm{odd}_\ell(G[t])}\le t(1+D+\dots+D^\ell)\delta N\le\delta_t\,tN$, and $tN\ge N\ge N_0$.

(c) Steps~2 and~4 of the proof of \cref{thm:A} do not use bipartiteness. Step~2, applied to the whole sequence, gives
$\limsup_n(1-\rank(M_n)/N_n)\le\mathbb{E}_\mu\langle P_K\delta_o,\delta_o\rangle$. Step~4, run in
$\mathrm{Alg}(\hat\mu)$ with $\iota^{-1}(1)$ in place of $V_+$ and $J=\iota^{-1}(0)$ (every neighbour of a vertex of
$\iota^{-1}(1)$ lies in $J$, since $\iota^{-1}(1)$ is independent), gives $\mathbb{E}_\mu\langle P_K\delta_o,\delta_o\rangle\le\hat\mu(\iota(o)=0)
=1-\hat\mu(\iota(o)=1)$. Take the supremum over $\hat\mu$.

(d) Ordering the vertices in matched pairs, $I+A_P$ is block diagonal with one block
$\bigl(\begin{smallmatrix}1&1\\1&1\end{smallmatrix}\bigr)$ of rank $1$ per edge of $P$. A regular bipartite graph with
at least one edge has a perfect matching, by Hall's theorem.
\end{informalproof}

\begin{remark}[On the extensions]\label{rem:ana-A-ext}
Part~(a) says that boundedness is needed only up to a diagonal rescaling and outside $o(N)$ rows. A rescaling changes
$M_{vu}$ into $s_v^{-1}M_{vu}s_u$, so it preserves the products of entries around closed walks, for instance
$M_{vu}M_{uv}$; ``bounded in some gauge'' is a condition on these products. In part~(b) the constants depend on $t$
through the degree $tD$ of $G[t]$, so the argument says nothing about the linear capacity, a supremum over~$t$, or
about codes whose number of symbols per vertex grows with $N$. Part~(c) explains why the limit object is needed. A
finite independent set $Y$ gives only $\rank M\ge\card Y$, since $M[Y,Y]$ is the identity, and graphs of large girth
can have small independence ratio: there are cubic graphs of arbitrarily large girth with independence ratio less than
$6/13$~\cite{Bollobas1981}. In a limit concentrated on the $d$-regular tree, by contrast, $\alpha_U=\frac12$: the
coin-flip colouring of Step~3 gives $\alpha_U\ge\frac12$, and the mass transport that sends one unit from each vertex
of $\iota^{-1}(1)$ to each of its $d$ neighbours gives $d\,\hat\mu(\iota(o)=1)\le d\,\hat\mu(\iota(o)=0)$. By part~(d)
the bound is attained along locally tree-like sequences as well, for instance on bipartite double covers of $d$-regular
graphs of large girth.
\end{remark}

\begin{remark}[Literature]\label{rem:ana-A-lit}
We have not found Theorem~A in the literature; its proof combines the easy half of
L\"uck approximation~\cite{Luck1994,BLS2011} with a tracial form of the Haemers bound~\cite{Haemers1979},
$\rank M\ge\card Y$ for an independent set $Y$, applied to an invariant random independent set of the limit; a tracial
version of the Haemers bound for finite graphs was studied by Gao, Gribling and Li~\cite{GaoGriblingLi2022}. The
positive semidefinite analogue is classical and exact: a positive semidefinite matrix with unit diagonal whose
off-diagonal support lies in the edges of a triangle-free graph has rank at least $N/2$ (Rosenfeld~\cite{Rosenfeld1991}
for real matrices, and Deaett~\cite{Deaett2011}). It needs no degree or coefficient bound, but it does not imply
Theorem~A, which concerns matrices that need not be symmetric or positive semidefinite. For such matrices
triangle-freeness alone does not suffice without a degree bound: Codenotti, Pudl\'ak and Resta~\cite{CPR2000} construct
triangle-free graphs for which $A-2I$, a matrix that fits the graph after scaling by $-\frac12$, has rank
$O(N^{3/4})$.
\end{remark}

\subsection{Codes over fields: torsion and odd characteristic}\label{sec:ana-fields}

We record what Theorem~A says about codes, and what it implies for codes over finite fields.

\begin{corollary}[Bounded real and complex linear codes]\label{cor:ana-codes}
\notformalized
Let $F\in\{\R,\C\}$, let $D,t\ge1$ and $B\ge1$ be fixed, and let $(G_n)$ be a locally bipartite sequence of graphs of
maximum degree at most $D$ on $N_n$ vertices. For $x\in F^{V\times[t]}$ write $x_v\in F^t$ for its block at $v$.
\begin{enumerate}[label=(\roman*)]
\item Let $C_n\subseteq F^{V(G_n)\times[t]}$ be a linear storage code with $t$ symbols per vertex and bounded
  coefficients: there are $t\times t$ matrices $R_{vu}$ ($uv\in E(G_n)$) with entries of modulus at most $B$ such that
  $x_v=\sum_{u\sim v}R_{vu}x_u$ for every $x\in C_n$ and every $v$. Then $\dim C_n\le(\frac12+o(1))\,tN_n$.
\item Let every vertex $v$ of $G_n$ be a receiver that wants $x_v$ and knows the blocks $x_u$ for $u$ in a set
  $I_n(v)\subseteq N_{G_n}(v)$ (directed side information). Let a linear map $E\colon F^{V(G_n)\times[t]}\to F^r$ be a
  linear index code with decoders $x_v=\Lambda_v(Ex)+\sum_{u\in I_n(v)}R_{vu}x_u$, valid for all $x$, where the maps
  $\Lambda_v$ are linear and the matrices $R_{vu}$ have entries of modulus at most $B$. Then $r\ge(\frac12-o(1))\,tN_n$.
\end{enumerate}
\end{corollary}

\begin{informalproof}
In both cases let $M_n$ be the matrix with blocks $(M_n)_{vv}=I_t$, $(M_n)_{vu}=-R_{vu}$ for the pairs $(v,u)$ that carry
a coefficient, and zero blocks elsewhere; it satisfies the hypotheses of \cref{prop:ana-A-ext}(b), so
$\rank M_n\ge(\frac12-o(1))\,tN_n$. (i) $C_n\subseteq\ker M_n$. (ii) For every $x$, the block of $M_nx$ at $v$ is
$\Lambda_v(Ex)$; so $M_n=\Lambda E$, where $\Lambda$ stacks the maps $\Lambda_v$, and $\rank M_n\le r$.
\end{informalproof}

By \cref{prop:ana-A-ext}(a), for $t=1$ the coefficients need to be bounded only after a diagonal rescaling and outside
$o(N_n)$ vertices. The statement does not cover coefficients that no rescaling makes bounded, numbers of symbols per
vertex that grow with $N_n$, codes over finite fields, or nonlinear codes. For finite fields Theorem~A still says
something: a gain must be a torsion effect.

\begin{corollary}[Torsion-only principle]\label{cor:torsion}
\notformalized
Fix a prime $p$, integers $D,t\ge1$, a real $B\ge1$ and $c>0$. Let $(G_n)$ be a locally bipartite sequence of graphs
of maximum degree at most $D$ on $N_n$ vertices. Let $M_n$ be matrices over $\F_p$ with rows and columns indexed by
$V(G_n)\times[t]$ whose $t\times t$ blocks satisfy $(M_n)_{vv}=I_t$ and $(M_n)_{vu}=0$ whenever $u\ne v$ and
$uv\notin E(G_n)$, and suppose that $\dim\ker M_n\ge(\frac12+c)\,tN_n$ for all $n$. Let $\tilde M_n$ be integer
matrices whose diagonal blocks are $I_t$, whose off-diagonal blocks vanish at non-edges and whose entries have modulus at
most $B$, that reduce to $M_n$ modulo $p$ (for instance the lifts with these diagonal blocks and entries in $[-p/2,p/2]$,
with $B=p/2$). Then $p$ divides at least $(c-o(1))\,tN_n$ of the invariant
factors of $\tilde M_n$; equivalently, the cokernel $\Z^{tN_n}/\tilde M_n\Z^{tN_n}$ contains a subgroup isomorphic to
$(\Z/p\Z)^{r_n}$ with $r_n\ge(c-o(1))\,tN_n$.
\end{corollary}

\begin{informalproof}
Write $\tilde M_n=U\Sigma W$ in Smith normal form, with $U,W$ invertible over $\Z$ and
$\Sigma=\mathrm{diag}(s_1,\dots,s_\rho,0,\dots,0)$, where $s_i\ge1$ and $\rho=\rank_\Q\tilde M_n$. Reducing modulo $p$,
where $U$ and $W$ stay invertible, gives $\rank_{\F_p}M_n=\card{\{i\le\rho:p\nmid s_i\}}$, so the number of invariant
factors divisible by $p$ is $\rho-\rank_{\F_p}M_n$. As a complex matrix, $\tilde M_n$ satisfies the hypotheses of
\cref{prop:ana-A-ext}(b), so $\rho\ge(\frac12-o(1))\,tN_n$, while $\rank_{\F_p}M_n=tN_n-\dim\ker M_n\le(\frac12-c)\,tN_n$.
The cokernel is isomorphic to $\Z^{tN_n-\rho}\oplus\bigoplus_i\Z/s_i\Z$, and its elements of order dividing $p$ form a
group $(\Z/p\Z)^{r_n}$ with $r_n=\card{\{i:p\mid s_i\}}$.
\end{informalproof}

An $\F_{p^e}$-linear storage code becomes, after the choice of a basis of $\F_{p^e}$ over $\F_p$, an $\F_p$-linear code
with $t=e$ symbols per vertex and the same rate, because multiplication by an element of $\F_{p^e}$ is $\F_p$-linear; so
\cref{cor:torsion} covers $\F_{p^e}$ for fixed $e$.

\begin{remark}[The Clebsch code is inherited from $\Q$]\label{rem:ana-clebsch}
The Clebsch graph is strongly regular with parameters $(16,5,0,2)$, so its adjacency matrix satisfies
$A^2=5I+2(J-I-A)=3I-2A+2J$, where $J$ is the all-ones matrix. Hence $(I+A)^2\equiv0\pmod 2$, and over $\F_2$ the image
of $I+A$ lies in its kernel. In fact $\rank_{\F_2}(I+A)=6$, which \cref{cor:H4} establishes formally; directly, the Clebsch graph is the Cayley
graph of
$\Z_2^4$ with connection set $\{e_1,\dots,e_4,e_1+\dots+e_4\}$, and in the group algebra
$\F_2[\Z_2^4]=\F_2[y_1,\dots,y_4]/(y_1^2,\dots,y_4^2)$, with $e_i\mapsto1+y_i$, the matrix $I+A$ acts as multiplication
by $a=\sigma_2+\sigma_3+\sigma_4$, where $\sigma_k$ is the $k$-th elementary symmetric polynomial in the $y_i$. The
product of $a$ with a monomial of degree $2$ is $y_1y_2y_3y_4$, and with a monomial of higher degree it is $0$, so the
image is spanned by $a$, $y_1a,\dots,y_4a$ and $y_1y_2y_3y_4$, which are linearly independent. The kernel of $I+A$, of
dimension $16-6=10>8$, is the code behind \cref{cor:H4}, and it is not a torsion effect. Over $\Q$ the eigenvalues of $A$ are $5$, $1$
and $-3$, with multiplicities $1$, $10$ and $5$, so $\rank_\Q(I-A)=6$. The integer kernel of $I-A$ is a direct summand
of $\Z^{16}$ of rank $10$, so its reduction modulo~$2$ is a $10$-dimensional subspace of $\ker_{\F_2}(I-A)=\ker_{\F_2}(I+A)$,
hence all of it. The binary code is therefore the reduction of the rational eigenvalue code of the eigenvalue~$1$
(\cref{prop:E}), whose recovery coefficients are $1$; \cref{cor:torsion} rules out exactly this mechanism on locally
bipartite families of bounded degree. The other lift $I+A$, with eigenvalues $6$, $2$ and $-2$, has full rank and ten
invariant factors divisible by~$2$: the torsion counted in \cref{cor:torsion} depends on the lift, and its conclusion
holds for every bounded lift. A single graph of girth~$4$ is not a locally bipartite family; the Clebsch graph shows
that local bipartiteness cannot be dropped from Theorem~A and \cref{cor:torsion}, even for bounded rational codes.
\end{remark}

The nilpotent identity $(I+W)^2=0$, which the Clebsch check matrix satisfies modulo~$2$, gives no gain once the girth is
at least~$5$.

\begin{lemma}[Index-two nilpotence at girth five]\label{lem:ana-matching}
\notformalized
Let $G$ be a graph of girth at least $5$, and let $W\in F^{V\times V}$, over any field $F$, have zero diagonal and
$W_{vu}=0$ whenever $uv\notin E(G)$. If $W^2=I$, then there is a perfect matching $P$ of $G$ such that $W_{vu}\ne0$
exactly when $uv\in P$, and $W_{vu}W_{uv}=1$ for $uv\in P$; in particular $\rank(I+W)=\rank(I-W)=N/2$. Over a field of
characteristic $2$, $(I+W)^2=I+W^2$, so every matrix $I+W$ that fits $G$ and satisfies $(I+W)^2=0$ has rank exactly
$N/2$.
\end{lemma}

\begin{informalproof}
For $v\ne w$, $(W^2)_{vw}$ is a sum over the common neighbours of $v$ and $w$, of which there is at most one, as $G$
has no cycle of length $3$ or $4$. Since $(W^2)_{vv}=\sum_{u\sim v}W_{vu}W_{uv}=1$, some neighbour $m(v)$ of $v$ has
$W_{v\,m(v)}W_{m(v)\,v}\ne0$. If $W_{vu'}\ne0$ for a neighbour $u'\ne m(v)$, then $m(v)$ and $u'$ are non-adjacent with
unique common neighbour $v$, and $(W^2)_{m(v)\,u'}=W_{m(v)\,v}W_{vu'}\ne0$, contradicting $W^2=I$; so row $v$ of $W$ is
supported on $\{m(v)\}$. Likewise, if $W_{u'v}\ne0$ for a neighbour $u'\ne m(v)$, then
$(W^2)_{u'\,m(v)}=W_{u'v}W_{v\,m(v)}\ne0$; so column $v$ is supported on $\{m(v)\}$. As $W_{m(v)\,v}\ne0$, row $m(v)$ is
supported on $\{v\}$, so $m(m(v))=v$, the pairs $\{v,m(v)\}$ form a perfect matching, and
$W_{v\,m(v)}W_{m(v)\,v}=(W^2)_{vv}=1$. Ordering the vertices in matched pairs, $I\pm W$ is block diagonal with blocks
$\bigl(\begin{smallmatrix}1&\pm a\\\pm a^{-1}&1\end{smallmatrix}\bigr)$ of rank $1$.
\end{informalproof}

In odd characteristic there is a different constraint, which compares a code with its sign twin.

\begin{proposition}[Sign twins in odd characteristic]\label{prop:oddchar}
\notformalized
Let $F$ be a field of characteristic different from $2$ and let $W\in F^{V\times V}$. Then
\[
  \rank(I+W)+\rank(I-W)\;=\;\rank\begin{pmatrix}I&W\\W&I\end{pmatrix}\;\ge\;N .
\]
Hence $\dim\ker(I+W)+\dim\ker(I-W)\le N$, and at most one of $\ker(I+W)$ and $\ker(I-W)$ has dimension larger than
$N/2$. If $W$ has zero diagonal and is supported on the edges of a graph $G$, the block matrix fits the bipartite double
cover $G\times K_2$, the graph on $V\times\{0,1\}$ in which $(v,i)\sim(u,j)$ if and only if $uv\in E(G)$ and $i\ne j$.
\end{proposition}

\begin{informalproof}
Let $Q=\bigl(\begin{smallmatrix}I&I\\I&-I\end{smallmatrix}\bigr)$; then $Q^2=2I$, so $Q$ is invertible. A direct
computation gives $Q\bigl(\begin{smallmatrix}I&W\\W&I\end{smallmatrix}\bigr)Q^{-1}
=\bigl(\begin{smallmatrix}I+W&0\\0&I-W\end{smallmatrix}\bigr)$, which proves the equality. The block matrix contains
$I$ as a principal $N\times N$ submatrix, so its rank is at least $N$. The last claim holds because the off-diagonal
entries of the block matrix are the entries $W_{vu}$ at positions $((v,0),(u,1))$ and $((v,1),(u,0))$.
\end{informalproof}

\begin{remark}[Gauge invariance]\label{rem:ana-gauge}
Over a field of characteristic different from~$2$, if $G$ is bipartite with a proper colouring $s\colon V\to\{\pm1\}$, the diagonal matrix $S=\mathrm{diag}(s)$ satisfies
$SWS=-W$ for $W$ supported on the edges, so $\ker(I+W)$ and $\ker(I-W)$ have the same dimension, at most $N/2$ by
\cref{prop:oddchar}. On a locally bipartite sequence the same conclusion up to $o(N_n)$, for all matrices over a field
of characteristic different from $2$ and uniformly in the number of symbols per vertex, would follow from approximate
gauge invariance: $\dim\ker(I+W_n)-\dim\ker(I-W_n)=o(N_n)$. In the bipartite limit object of the proof of \cref{thm:A}
the coin-flip colouring does conjugate $I+W$ into $I-W$, although that proof uses the colouring differently, through the
independent set $\beta^{-1}(+1)$. In characteristic $2$,
$I+W=I-W$ and \cref{prop:oddchar} says nothing; on the Clebsch graph $I+A=I-A$ has a kernel of dimension $10>8$ over
$\F_2$.
\end{remark}

\paragraph{Why it matters.} For constructions, \cref{cor:torsion} is the finite-field part of the design rule stated
after \cref{thm:A}; in particular, a gain cannot be inherited by reduction modulo~$p$ from a code over $\Q$ with bounded
coefficients, as the Clebsch code is (\cref{rem:ana-clebsch}). In characteristic~$2$, a gain also cannot come from the identity
$(I+W)^2=0$ once the girth is at least~$5$ (\cref{lem:ana-matching}). For the impossibility side,
\cref{prop:oddchar} shows that in odd characteristic a gain on one sign twin is paid for on the other. In
characteristic~$2$ only \cref{cor:torsion} applies.

\subsection{Limits of local and spectral arguments}\label{sec:ana-limits}

Theorem~S settles all scalar linear codes on triangle-free graphs of maximum degree~$3$, and Theorem~A settles bounded
complex linear codes on locally bipartite graphs of every bounded degree by a spectral argument. The undirected crux needs finite
fields (indeed all codes), at degree at least~$4$ for scalar linear codes and already at degree~$3$ for vector-linear
and nonlinear ones. This subsection
records three limits: inside tree-like regions, the local necessary conditions of a smallest counterexample are
compatible with rate above $\frac12+o(1)$ at degree at least~$4$ (\cref{prop:L}); real eigenvalue codes, the mechanism
behind several known codes that beat $\frac12$ at bounded degree, cannot reach rate $\frac12$ at girth~$7$ or more
(\cref{prop:E}); and girth~$5$ does not replace low degree (\cref{prop:HS}).

A class $\mathcal K$ of pairs $(G,M)$, where $M$ fits $G$ over some field, is \emph{hereditary} if
$(G[W],M[W,W])\in\mathcal K$ whenever $(G,M)\in\mathcal K$ and $W\subseteq V(G)$ is nonempty; for example, triangle-free
graphs of maximum degree at most $4$ with all matrices that fit them, over all fields. For $C=\ker M$ and
$U\subseteq V$ let $r(U)=\dim C|_U$ be the dimension of the restriction of $C$ to the coordinates in $U$ (the rank
function of the matroid of $C$), and let $C^\perp=\{y\in F^V:\sum_vy_vx_v=0\text{ for all }x\in C\}$, the row space of
$M$. For a hereditary class, the floor bound $\dim\ker M\le\lfloor N/2\rfloor$ for all
members implies $\dim\ker M\le\nu(G)$, the matching number, by the argument of \cref{thm:S-barrier}: for every vertex
set $U$, $\dim\ker M\le\card U+\sum_H\lfloor\card{H}/2\rfloor$ over the components $H$ of $G-U$, and the Tutte--Berge
formula turns the minimum over $U$ into $\nu(G)$. The floor bound is therefore the statement to prove.

\begin{lemma}[Smallest counterexamples]\label{lem:ana-mincex}
\notformalized
Let $\mathcal K$ be hereditary, and let $(G,M)\in\mathcal K$ satisfy $\dim\ker M>\lfloor N/2\rfloor$ with $N$ as small as
possible. Then
\begin{enumerate}[label=(\alph*)]
\item $N$ is odd and $\dim\ker M=(N+1)/2$;
\item $r(U)\ge\lfloor\card U/2\rfloor+1$ for every nonempty $U\subseteq V$;
\item for every nonempty $X\subseteq V$, the vectors of $C^\perp$ supported in $X$, in particular the rows of $M$
  supported in $X$, span a space of dimension at most $\lceil\card X/2\rceil-1$.
\end{enumerate}
\end{lemma}

\begin{informalproof}
A $1\times1$ matrix that fits a graph is $(1)$, so $N\ge2$. For $U\subseteq V$, the codewords vanishing on $U$ form a
subspace of dimension $\dim C-r(U)$, and their restrictions to $V\setminus U$ lie in $\ker M[V\setminus U,V\setminus U]$
and determine them; so for $U\ne V$ the pair $(G-U,M[V\setminus U,V\setminus U])\in\mathcal K$ has a kernel of dimension
at least $\dim C-r(U)$. (a) If $N$ were even, deleting one vertex ($r\le1$) would leave a kernel of dimension at least
$N/2>\lfloor(N-1)/2\rfloor$ on fewer vertices; if $\dim C\ge(N+3)/2$, deleting one vertex would leave a kernel of
dimension at least $(N+1)/2>\lfloor(N-1)/2\rfloor$. (b) By minimality for $U\ne V$, and trivially for $U=V$,
$\dim C-r(U)\le\lfloor(N-\card U)/2\rfloor$, so $r(U)\ge(N+1)/2-\lfloor(N-\card U)/2\rfloor=\lfloor\card U/2\rfloor+1$,
as $N$ is odd. (c) The vectors of $C^\perp$ supported in $X$ form the annihilator of the
restriction $C|_X$ in $F^X$, of dimension $\card X-r(X)\le\card X-\lfloor\card X/2\rfloor-1$.
\end{informalproof}

One way to reach a contradiction in a smallest counterexample is to exhibit a set $X$ with $\lceil\card X/2\rceil$
independent vectors of $C^\perp$ supported in $X$, for instance rows of $M$. The proof of Theorem~S does this with a
maximal chain grown from a single vertex along recovery neighbourhoods; such a chain closes only through cycles or the
boundary of the vertex set, never inside a tree-like region. The next proposition shows that inside a tree-like region of degree
$d\ge4$ no such set exists, however the region is explored, and that the local conditions there are met by a code of
rate close to $(d-2)/(d-1)$.

\begin{proposition}[Proposition L: tree codes are locally consistent]\label{prop:L}
\notformalized
Let $d\ge3$ and $r\ge1$, let $B$ be the ball of radius $r$ around the root of the $d$-regular tree, $B^\circ\subseteq B$
the set of vertices at distance less than $r$ from the root, and $F$ a field. For each $v\in B^\circ$ let $z_v\in F^B$
have $z_v(v)=1$, $z_v(u)\ne0$ for every neighbour $u$ of $v$, and $z_v(w)=0$ otherwise (over $\F_2$ these are the rows
of $I+A$ indexed by $B^\circ$). Let $C_B=\{x\in F^B:\sum_wz_v(w)x_w=0\text{ for all }v\in B^\circ\}$ and
$r(U)=\dim C_B|_U$.
\begin{enumerate}[label=(\roman*)]
\item $\dim C_B=d(d-1)^{r-1}$, the number of vertices at distance $r$, and $\dim C_B/\card B\to(d-2)/(d-1)$ as
  $r\to\infty$: this is $\frac12$ for $d=3$ and $\frac23$ for $d=4$.
\item $r(U)\ge\card U-\lfloor\card U/(d-1)\rfloor$ for every $U\subseteq B$.
\item If $d\ge4$, then $r(U)\ge\lfloor\card U/2\rfloor+1$ for every nonempty $U\subseteq B$, and for every nonempty
  $X\subseteq B$ the vectors of $C_B^\perp$ supported in $X$, in particular the vectors $z_v$ supported in $X$, span a
  space of dimension at most $\lceil\card X/2\rceil-1$.
\end{enumerate}
\end{proposition}

\begin{informalproof}
Orient the tree away from the root; for $v\in B^\circ$ let $\mathrm{ch}(v)$ be its set of children, of size $d$ at the
root and $d-1$ elsewhere. (i) Order $B^\circ$ by distance from the root. A child of $v$ lies in the support of $z_w$, for
$w\in B^\circ$, only if $w=v$, or $w$ is that child or one of its children, which come after $v$; so each $z_v$ is
nonzero at a coordinate where all earlier vectors vanish, the vectors $z_v$ are linearly independent, and
$\dim C_B=\card B-\card{B^\circ}$ is the number of vertices at distance $r$. Since $\card B=1+d((d-1)^r-1)/(d-2)$, the
ratio tends to $(d-2)/(d-1)$.

(ii) The annihilator of $C_B|_U$ in $F^U$ is $C_B^\perp\cap F^U$, where $C_B^\perp$ is spanned by the vectors $z_v$, so
$r(U)=\card U-\dim(C_B^\perp\cap F^U)$. Let $D_U=\{v\in B^\circ:\mathrm{ch}(v)\subseteq U\}$. If
$y=\sum_v\lambda_vz_v$ is supported in $U$ and $\lambda\ne0$, let $v$ be a vertex of maximal depth with $\lambda_v\ne0$.
For a child $c$ of $v$, the only vectors $z_w$ with $z_w(c)\ne0$ are $z_v$, $z_c$ and the $z_w$ for the children $w$ of
$c$, and $\lambda_c=\lambda_w=0$ because $c$ and $w$ are deeper than $v$; so $y_c=\lambda_vz_v(c)\ne0$, $c\in U$, and
$v\in D_U$. Since the vectors $z_v$ are independent by (i), each element of $C_B^\perp\cap F^U$ has a unique coefficient
vector $\lambda$, and $\lambda\mapsto\lambda|_{D_U}$ is injective on these coefficient vectors; so
$\dim(C_B^\perp\cap F^U)\le\card{D_U}$. The children sets of distinct vertices are disjoint subsets of size at least
$d-1$, so $\card{D_U}\le\lfloor\card U/(d-1)\rfloor$.

(iii) For $d\ge4$, (ii) gives $r(U)\ge\card U-\lfloor\card U/3\rfloor$, which is at least $\lfloor\card U/2\rfloor+1$:
check $\card U\le5$ directly, and for $\card U\ge6$ use $\card U-\lfloor\card U/3\rfloor\ge2\card U/3\ge\card U/2+1$.
The second claim follows as in \cref{lem:ana-mincex}(c).
\end{informalproof}

\paragraph{Why it matters.} In a $d$-regular graph of girth at least $2r+2$, every ball of radius $r$ is a copy of $B$.
By \cref{prop:L}, inside such a ball the constraints that a smallest counterexample satisfies, namely (b) and (c) of
\cref{lem:ana-mincex} for subsets of the ball and the recoverability of every interior vertex from its neighbours, are
all met by a code of rate close to $(d-2)/(d-1)$, which exceeds $\frac12$ when $d\ge4$. So an argument that uses only
these local necessary conditions inside tree-like regions cannot give even the asymptotic bound $\frac12+o(1)$ at degree
at least~$4$. At degree~$3$ the ball code $C_B$ has rate $\frac12+1/\card B$, so such conditions are compatible with the
asymptotic bound there; the exact floor bound needs global information even at degree~$3$, as the closing of the maximal
chain in the proof of Theorem~S shows. For such arguments, large girth is the hardest case, not the easiest. And the
degree-$4$ statement over $\F_2$ on $4$-regular graphs of large girth is the scalar-linear, characteristic-$2$ case of the
undirected crux at degree~$4$, for which no characteristic-$p$ analogue of Theorem~A is known.

\begin{remark}[Computational evidence at degree four]\label{rem:ana-deg4}
Exhaustive exact searches with SAT solvers, not independently verified and not used in any proof, found no
counterexample to the floor bound for triangle-free graphs of maximum degree $4$: over $\F_2$ none of the
$193{,}900$ connected $4$-regular triangle-free graphs on $17$ vertices has a linear storage code of dimension $9$; over
$\F_3$ none of the connected $4$-regular triangle-free graphs on $13$ and $15$ vertices, and none of the non-bipartite
connected triangle-free graphs with degrees between $2$ and $4$ on $9$ vertices, or on $11$ vertices with a vertex of
degree $4$, has a code of dimension larger than its matching number.
\end{remark}

Several known codes that beat $\frac12$ on graphs of bounded degree, for instance on the Clebsch graph over $\Q$ and on
the Hoffman--Singleton graph, come from eigenvalues. For a graph $G$ with adjacency matrix $A$ and a real number
$\theta\ne0$, the matrix $I-A/\theta$ fits $G$, has coefficients of modulus $1/\abs\theta$, and its kernel, the \emph{eigenvalue code} of $\theta$, is the
$\theta$-eigenspace of $A$, of dimension $\mathrm{mult}(\theta)$ (zero if $\theta$ is not an eigenvalue). The next
proposition, a linear-programming bound in the spirit of Nozaki's bounds for regular graphs~\cite{Nozaki2015}, shows
where this mechanism stops. Let
\[
  p_0=1,\quad p_1(x)=\frac{x}{\sqrt d},\quad p_2(x)=\frac{x^2-d}{\sqrt{d(d-1)}},\quad
  p_3(x)=\frac{x(x^2-2d+1)}{\sqrt d\,(d-1)},\qquad K_n(\theta)=\sum_{k=0}^np_k(\theta)^2 ;
\]
these are the first orthonormal polynomials of the Kesten--McKay law, the spectral measure of the $d$-regular tree at a
vertex~\cite{Kesten1959,McKay1981}.

\begin{proposition}[Eigenvalue codes and girth]\label{prop:E}
\notformalized
Let $d\ge3$ and let $G$ be a $d$-regular graph on $N$ vertices with girth at least $g$.
\begin{enumerate}[label=(\roman*)]
\item For $n\in\{1,2,3\}$ with $2n+1\le g$ and every real $\theta$, $\mathrm{mult}(\theta)\le N/K_n(\theta)$.
\item $K_1(\theta)<2$ if and only if $\theta^2<d$; and $K_2(\theta)<2$ if and only if $1<\theta^2<d$.
\item $K_3(\theta)>2$ for every real $\theta$. Hence, if $g\ge7$, every eigenvalue has multiplicity less than $N/2$.
\item If $g\ge7$, no eigenvalue code of $G$ has dimension $N/2$ or more. If $g\ge5$ and an eigenvalue code has dimension
  larger than $N/2$, then $\theta$ is an integer with $2\le\abs\theta<\sqrt d$, so $d\ge5$.
\end{enumerate}
Equality holds in (i) for the Petersen graph ($d=3$, $g=5$, $n=2$, $\theta=1$, multiplicity $5$) and for the
Hoffman--Singleton graph ($d=7$, $g=5$, $n=2$, $\theta=2$, multiplicity $28$).
\end{proposition}

\begin{informalproof}
\emph{Orthonormality.} Let $A_T$ be the adjacency operator of the $d$-regular tree, $o$ its root, $S_k$ the sphere of
radius $k$ around $o$, and $e_k=\card{S_k}^{-1/2}\mathbf 1_{S_k}$. Counting neighbours in adjacent spheres gives
$A_Te_0=\sqrt d\,e_1$, $A_Te_1=\sqrt d\,e_0+\sqrt{d-1}\,e_2$ and $A_Te_k=\sqrt{d-1}\,(e_{k-1}+e_{k+1})$ for $k\ge2$. A
direct computation from these relations gives $p_k(A_T)e_0=e_k$ for $k\le3$; for instance
$(A_T^2-d)e_0=\sqrt{d(d-1)}\,e_2$. Hence $\langle p_j(A_T)e_0,p_k(A_T)e_0\rangle=\delta_{jk}$ for $j,k\le3$.

\emph{Moments.} For $k<g$ and every vertex $v$ of $G$, $(A^k)_{vv}=\langle A_T^ke_0,e_0\rangle$. Indeed, let
$\pi\colon T\to G$ be a covering of the component of $v$ by the $d$-regular tree with $\pi(o)=v$. Walks from $v$ lift
uniquely to walks from $o$ of the same length, and a lift of a closed walk of length $k<g$ is closed: otherwise the path
in $T$ between its distinct endpoints would project to a closed non-backtracking walk of positive length at most $k$,
which contains a cycle of length at most $k<g$. Conversely, closed walks from $o$ project to closed walks from $v$.

(i) Let $q=\sum_{k\le n}c_kp_k$ with $c_k=p_k(\theta)/K_n(\theta)$, so that $q(\theta)=1$. The symmetric matrix $q(A)^2$
has eigenvalue $1$ on the $\theta$-eigenspace and nonnegative eigenvalues elsewhere, so
$\mathrm{mult}(\theta)\le\operatorname{tr}q(A)^2$. Since $q^2$ has degree $2n<g$, the moments give
$\operatorname{tr}q(A)^2=N\lVert q(A_T)e_0\rVert^2=N\sum_kc_k^2=N/K_n(\theta)$.

(ii) $K_1(\theta)=1+\theta^2/d$. With $s=\theta^2$, $d(d-1)\bigl(K_2(\theta)-2\bigr)=(d-1)s+(s-d)^2-d(d-1)=(s-1)(s-d)$.

(iii) With $s=\theta^2$, $d(d-1)^2\bigl(K_3(\theta)-2\bigr)=f(s)=s^3-(3d-1)s^2+(3d^2-4d+2)s+d(d-1)$, and
$f(d+y)=y^3+y^2-2(d-1)y+d(d-1)^2$. The derivative $3y^2+2y-2(d-1)$ has roots $y_\pm=(-1\pm\sqrt{6d-5})/3$, a local
maximum at $y_-<0$ and a local minimum at $y_+>0$. On $[-d,0]$ the minimum of $f(d+y)$ is therefore attained at an
endpoint, where the values are $f(0)=d(d-1)>0$ and $f(d)=d(d-1)^2>0$. On $[0,\infty)$ the minimum is
$f(d+y_+)=d(d-1)^2-y_+^2(2y_++1)$, using $3y_+^2+2y_+=2(d-1)$. From $y_+<\sqrt{6d}/3$ we get
$y_+^2(2y_++1)<\frac{2d}3\bigl(\frac{2\sqrt{6d}}3+1\bigr)$, which is less than $d(d-1)^2$ because
$(d-1)^2>\frac49\sqrt{6d}+\frac23$ for $d\ge3$ (at $d=3$, $4>2.56$, and the left side grows faster). So $f(s)>0$ for
all $s\ge0$. The second claim follows from (i) with $n=3$.

(iv) The first claim is (iii). For the second: $A$ has integer entries, so the roots of the minimal polynomial of
$\theta$ over $\Q$ are eigenvalues of the same multiplicity; if this multiplicity exceeds $N/2$, the minimal polynomial
has degree $1$, and $\theta$, an algebraic integer, is an integer. By (i) with $n=2$ and (ii), $1<\theta^2<d$, so
$\abs\theta\ge2$ and $d>4$.

\emph{Equality cases.} $K_2(1)=1+\frac13+\frac46=2$ for $d=3$, and $K_2(2)=1+\frac47+\frac9{42}=\frac{25}{14}$ for
$d=7$, with $50\cdot\frac{14}{25}=28$; the multiplicities are those of the Petersen graph (spectrum $3$, $1^5$,
$(-2)^4$) and of the Hoffman--Singleton graph (computed in the proof of \cref{prop:HS}).
\end{informalproof}

\paragraph{Why it matters.} \Cref{prop:E} bounds the girth up to which the spectral mechanism can work: real eigenvalue codes,
which give the Clebsch graph a code of dimension $10$ over $\Q$ (eigenvalue $1$) and the Hoffman--Singleton graph one of
dimension $28$ (\cref{prop:HS}), never reach rate $\frac12$ at girth $7$ or more, at any degree, and beat it at girth $5$
or $6$ only when $d\ge5$; whether they beat it at girth~$6$ at all is open. It is a finite-girth counterpart of
Theorem~A for eigenvalue codes, with explicit thresholds.
It does not see $p$-torsion: over $\F_p$ the rank of $A-\theta I$ can be smaller than over $\Q$, and over $\F_2$ a code
need not come from an eigenvalue at all, and by \cref{cor:torsion} the finite-field gains that matter here are torsion
effects.

\begin{proposition}[Proposition HS: Hoffman--Singleton beats one half outside characteristic two]\label{prop:HS}
\notformalized
Let $A$ be the adjacency matrix of the Hoffman--Singleton graph ($50$ vertices, $7$-regular, girth $5$). Over every field
$F$ of characteristic different from $2$, the matrix $I-A/2$ fits the graph and has rank at most $22$, with equality if
the characteristic is $0$. So $\ker(I-A/2)$ is a linear storage code over $F$ of dimension at least $28>25$, with
recovery coefficients $\frac12$, and disjoint unions of copies of the graph form a $7$-regular family of girth $5$ that
carries such codes of rate at least $0.56$.
\end{proposition}

\begin{informalproof}
By~\cite{HoffmanSingleton1960} the Hoffman--Singleton graph is strongly regular with parameters $(50,7,0,1)$:
adjacent vertices have no common neighbour and non-adjacent vertices exactly one. Hence $A^2=7I+(J-I-A)$, where $J$ is
the all-ones matrix, that is, $(A-2I)(A+3I)=J$. The all-ones vector is an eigenvector with eigenvalue $7$, and $J$
vanishes on its orthogonal complement, which is invariant under $A$; so the other eigenvalues are $2$ and $-3$, with
multiplicities $f+g=49$ and $7+2f-3g=\operatorname{tr}A=0$, and $f=28$. Thus
$\rank_\Q(2I-A)=22$. Minors of an integer matrix that vanish over $\Z$ vanish modulo $p$, so
$\rank_{\F_p}(2I-A)\le22$, and $I-A/2=(2I-A)/2$ has the same rank whenever $2$ is invertible. The rank over a field
depends only on its characteristic.
\end{informalproof}

\begin{remark}[Computations on the Hoffman--Singleton graph]\label{rem:ana-hs}
By exact computation (not used in any proof), the rank of $2I-A$ is $22$ over $\F_3$ and $\F_7$ but $21$ over $\F_5$,
an example of the $p$-torsion that \cref{prop:E} does not see. Over $\F_2$ the matrix $I+A$ has rank $29$, giving only
dimension $21$, and we do not know whether the Hoffman--Singleton graph carries an $\F_2$-linear storage code of
dimension larger than $25$.
\end{remark}

\paragraph{Why it matters.} Theorem~S needs maximum degree~$3$ (and no triangles), and Theorem~A needs local
bipartiteness. One might hope
that large girth alone gives the bound $\frac12$ at every bounded degree, over every field. \Cref{prop:HS} shows that
girth~$5$ does not suffice: at degree~$7$, over $\Q$ and over every field of odd characteristic, there are linear codes
of rate $0.56$ with bounded coefficients. By \cref{prop:E}, eigenvalue codes need girth at most~$6$ to beat $\frac12$;
whether, for each $d$, some girth forces rate at most $\frac12+o(1)$ for all linear codes over a fixed finite field is
open, and in characteristic~$2$ even the Hoffman--Singleton graph is undecided.

Taken together, on side information of bounded degree $d\ge4$ that is locally bipartite, an $\F_p$-linear code with a
fixed number of symbols per vertex that beats $\frac12$ by a constant must have $p$-torsion of linear rank in every
bounded integer lift of its check matrix (\cref{cor:torsion}); in characteristic~$2$ it gains nothing from the identity
$(I+W)^2=0$ once the girth is at least $5$ (\cref{lem:ana-matching}); it does not come from a real eigenvalue code once
the girth is at least $7$ (\cref{prop:E}); and the local necessary conditions of a smallest counterexample cannot refute
its existence inside tree-like regions (\cref{prop:L}). These are design
constraints for constructions, and they say where a characteristic-$p$ proof of the crux would have to work.

\subsection{Targets, Mycielskians and Question G}\label{sec:ana-targets}

By Theorem~H (\cref{thm:H}), a homomorphism from side information to a good graph $T$ certifies the rate bound $\frac12$
for every code on the side information, linear or not. If the homomorphism is defined only after the deletion of a
vertex set $U$, conditioning on the symbols at $U$ weakens the bound by at most $\card U/2$ symbols: for a valid
polymatroid~$h$, the function $S\mapsto h(S\cup U)-h(U)$ on subsets of $V\setminus U$ is a valid polymatroid for $G-U$,
so $h(V)\le h(U)+\frac12\sum_{v\notin U}h(\{v\})\le\frac12\sum_vh(\{v\})+\frac12\sum_{u\in U}h(\{u\})$. A single good target
that receives the side information of the crux after the deletion of $o(N)$ vertices would therefore prove the crux for
all codes. Two facts constrain the targets. Good graphs are triangle-free: the triangle carries the code
$\{(a,b,a+b)\}$ of rate $\frac23$, and a graph that contains a triangle receives it (\cref{thm:H}). And homomorphisms do not decrease
the fractional chromatic number $\chi_f$, the least total weight of a family of weighted independent sets that covers
every vertex with weight at least $1$, because preimages of independent sets are independent. So side information $H$
with a deleted set $U$ needs a target with $\chi_f(T)\ge\chi_f(H-U)\ge(\card{V(H)}-\card U)/\alpha(H)$, where $\alpha$
is the independence number.

The classical way to raise $\chi_f$ while keeping a graph triangle-free is Mycielski's construction. For a graph $G$ on
$V$, the \emph{Mycielskian} $\mathrm M(G)$ has vertex set $V\cup V'\cup\{z\}$, where $V'=\{v':v\in V\}$, and edges $uv$,
$uv'$ and $u'v$ for $uv\in E(G)$, and $v'z$ for $v\in V$. Thus $\mathrm M(K_2)$ is the $5$-cycle and $\mathrm M^2(K_2)$
is the Gr\"otzsch graph. The Mycielskian of a triangle-free graph is triangle-free~\cite{Mycielski1955}, and
$\chi_f(\mathrm M(G))=\chi_f(G)+1/\chi_f(G)$~\cite{LPU1995}, so $\chi_f(\mathrm M^k(K_2))\to\infty$.

\begin{lemma}[Homomorphisms into a Mycielskian]\label{lem:ana-myc}
\notformalized
For graphs $H$ and $G$, $H$ admits a homomorphism to $\mathrm M(G)$ if and only if $H$ has an independent set $Z$ such that
the set $N_H(Z)$ of vertices outside $Z$ with a neighbour in $Z$ is independent and $H-Z$ admits a homomorphism to $G$.
\end{lemma}

The `if' direction is the case $t=1$ of~\cite[Lemma~2.3]{GHW2025}; we include the short proof of both directions.

\begin{informalproof}
Let $\varphi\colon H\to\mathrm M(G)$, and put $Z=\varphi^{-1}(z)$ and $Y=\varphi^{-1}(V')$. Both are independent,
because $z$ has no loop and $V'$ is independent in $\mathrm M(G)$, and $N_H(Z)\subseteq Y$ because $z$ is adjacent only
to $V'$. Define $\psi$ on $H-Z$ by $\psi(w)=\varphi(w)$ if $\varphi(w)\in V$ and $\psi(w)=v$ if $\varphi(w)=v'$. An edge
of $H-Z$ has at most one end in $Y$; if it has none, $\varphi$ maps it to an edge of $G$; if it has one, $\varphi$ maps it
to an edge $uv'$ with $uv\in E(G)$, and $\psi$ maps it to $uv$. So $\psi$ is a homomorphism to $G$. Conversely, given
$Z$ and $\psi\colon H-Z\to G$, map $Z$ to $z$, each $w\in N_H(Z)$ to $\psi(w)'$, and every other vertex $w$ to $\psi(w)$.
There are no edges inside $Z$ or inside $N_H(Z)$, and none between $Z$ and the vertices outside $Z\cup N_H(Z)$; edges
between $Z$ and $N_H(Z)$ go to edges $zv'$; and every other edge $ww''$ goes to $\psi(w)\psi(w'')$, $\psi(w)'\psi(w'')$ or
$\psi(w)\psi(w'')'$, which are edges of $\mathrm M(G)$ because $\psi(w)\psi(w'')\in E(G)$.
\end{informalproof}

\begin{corollary}[Odd cycle transversals]\label{cor:ana-oct}
\notformalized
If $H$ is triangle-free and $H-X$ is bipartite for a set $X$ of $k$ vertices, then $H$ admits a homomorphism to
$\mathrm M^k(K_2)$.
\end{corollary}

\begin{informalproof}
Induction on $k$; a bipartite graph maps to $K_2$. For $k\ge1$ pick $x\in X$. The set $Z=\{x\}$ is independent, $N_H(x)$
is independent because $H$ is triangle-free, and $H-x$ is triangle-free and becomes bipartite after the deletion of
$X\setminus\{x\}$, so $H-x\to\mathrm M^{k-1}(K_2)$ by induction. \Cref{lem:ana-myc} gives $H\to\mathrm M^k(K_2)$. (A
similar iteration is used in~\cite{GHW2025}.)
\end{informalproof}

\begin{proposition}[Mycielskians leave the good graphs]\label{prop:M}
\notformalized
\begin{enumerate}[label=(\roman*)]
\item The Clebsch graph admits a homomorphism to $\mathrm M^6(K_2)$. Hence $\mathrm M^6(K_2)$ is not good.
\item $K_2$ and $\mathrm M(K_2)$, the $5$-cycle, are good. Hence there is $j\in\{1,\dots,5\}$ such that $\mathrm M^j(K_2)$
  is good and $\mathrm M^{j+1}(K_2)$ is not: the Mycielski construction does not preserve goodness.
\item No good graph receives every triangle-free graph of maximum degree $5$. More generally, if $T_1,T_2,\dots$ are
  graphs and every finite triangle-free graph admits a homomorphism to some $T_k$, then some $T_k$ is not good; by
  \cref{cor:ana-oct} this applies to the graphs $\mathrm M^k(K_2)$, which are universal for triangle-free graphs by
  a theorem of Cropper, Gyárfás and Lehel~\cite{CGL2006}.
\end{enumerate}
\end{proposition}

\begin{informalproof}
(i) The Clebsch graph $\Gamma$ is triangle-free and $5$-regular on $16$ vertices. For an edge $uv$, the neighbourhoods
$N(u)$ and $N(v)$ are disjoint independent sets (as $\Gamma$ is triangle-free), so $\Gamma[N(u)\cup N(v)]$ is bipartite
on $10$ vertices and the remaining set $X$ of $6$ vertices has $\Gamma-X$ bipartite. By \cref{cor:ana-oct},
$\Gamma\to\mathrm M^6(K_2)$. The Clebsch graph is not good: by \cref{cor:H4} it has guessing number $10>8$, while a good
graph on $16$ vertices has guessing number at most $8$ by \cref{thm:entropic-rate}. By \cref{thm:H}, a graph that
receives a graph that is not good is not good. (ii) The $5$-cycle is good (\cref{prop:C5}), and $K_2$ maps to it, so
$K_2$ is good by \cref{thm:H}. With (i), the least $k$ for which $\mathrm M^k(K_2)$ is not good satisfies $2\le k\le6$;
take $j=k-1$. (iii) The Clebsch graph is triangle-free of maximum degree $5$, and a graph that receives it is not good.
\end{informalproof}

\begin{remark}[Which iterate fails; computational]\label{rem:ana-myc}
The Gr\"otzsch graph $\mathrm M^2(K_2)$ is good by an exact certificate, a nonnegative rational combination of
elemental Shannon inequalities found by computer and checked in exact arithmetic; the certificate is not published here,
not formalized and not used in any proof. With it,
$j\in\{2,3,4,5\}$ in \cref{prop:M}(ii). A SAT computation (not used in any proof) indicates that the Clebsch graph does
not map to $\mathrm M^5(K_2)$, so the argument of (i) cannot locate the first iterate that is not good; which iterate it
is remains open. At the level of Shannon inequalities the construction may fail at the first step: numerically, by a
floating-point linear program without an exact certificate, the Mycielskian of the circular clique $K_{8/3}$ (which is
good by \cref{prop:targets}) has Shannon bound about $7/13>\frac12$.
\end{remark}

\paragraph{Why it matters.} Mycielski's construction is the standard way to build triangle-free graphs of large
fractional chromatic number, and its iterates are universal for triangle-free graphs~\cite{CGL2006}; \cref{cor:ana-oct} refines this, bounding
the number of iterations by the size of an odd cycle transversal.
\Cref{prop:M} shows that this universality forces them out of the good graphs. Targets of growing $\chi_f$ for the
homomorphism route to the crux must therefore be built differently, and they must avoid receiving the Clebsch graph and
the other graphs known not to be good.

\begin{remark}[Good graphs with $\chi_f>3$; computational]\label{rem:ana-circulants}
Write $C_n(S)$ for the Cayley graph of $\Z_n$ with connection set $\pm S$; for vertex-transitive graphs
$\chi_f=N/\alpha$. Exact rational certificates, nonnegative combinations of elemental Shannon inequalities found by
computer and checked in exact arithmetic, show that the triangle-free circulants
$C_{13}(1,5)$ ($4$-regular, $\chi_f=13/4$), $C_{16}(1,3,8)$ ($5$-regular, $\chi_f=16/5$) and $C_{20}(1,6,10)$
($5$-regular, $\alpha=6$, $\chi_f=10/3$) are good. The certificate for $C_{20}(1,6,10)$ consists of $777$ inequalities;
$10/3$ is the largest value of $\chi_f$ among good graphs known to us. The certificates are not published here, and none
of this is formalized or used in any proof. These examples show that $\chi_f\le3$ is no barrier, and that $\chi_f$ does not
determine goodness: $C_{16}(1,3,8)$ has the same number of vertices, degree, independence number and fractional
chromatic number as the Clebsch graph, which is not good. Nearby circulants are not good: for instance
$C_{20}(2,5,6)$, with $\chi_f=10/3$, carries the $\F_2$-linear code $\ker(I+A)$ of dimension $12>10$.
\end{remark}

\begin{remark}[Cayley graphs of elementary abelian $2$-groups; partly computational]\label{rem:ana-z2}
Proof sketch only: every non-bipartite Cayley graph $\Gamma=\mathrm{Cay}(\Z_2^k,S)$ is not good. The graph $\Gamma$ is
bipartite if and only if some linear functional equals $1$ on $S$, and by linear duality over $\F_2$ this fails exactly
when some subset of $S$ with an odd number of elements sums to $0$. A minimal such subset $S'$ consists of $m-1$ linearly independent
elements and their sum, with $m$ odd (and $m\ge5$ if $\Gamma$ is triangle-free), so $\mathrm{Cay}(\langle
S'\rangle,S')$ is the folded $m$-cube, a subgraph of $\Gamma$. In the group algebra the matrix $I+A$ of a folded cube
satisfies $(I+A)^2=0$ over $\F_2$, so its kernel has dimension at least half; exact computation shows that the dimension
is $(2^{m-1}+2^{(m-1)/2})/2$ for $m\le11$, and we have not proved this formula in general. Where it holds, the folded
$m$-cube is not good, and neither is $\Gamma$, by \cref{thm:H}. The argument is complete for $k\le10$, where $m\le11$;
it would exclude, as targets, all non-bipartite triangle-free Cayley graphs of elementary abelian $2$-groups, a family that contains
the Clebsch graph and the folded cubes.
\end{remark}

These examples lead to the following question.

\begin{question}[Question G: the fractional chromatic number of good graphs]\label{prob:G}
Is $\chi_f$ unbounded on good graphs? That is, is $\sup\{\chi_f(T):T\text{ a good graph}\}$ infinite?
\end{question}

Good graphs are triangle-free, and triangle-free graphs have unbounded $\chi_f$; the question is whether goodness bounds
it. The analogous question for digraphs, with the fractional dichromatic number (fractional covers by vertex sets
inducing acyclic subgraphs), is also open.

\paragraph{Why it matters.} A negative answer, $\chi_f(T)\le C$ for every good graph $T$, would rule out the homomorphism
route for every family of side information whose underlying graphs keep fractional chromatic number above $C$ after the
deletion of $o(N)$ vertices, for instance families whose independence ratio stays below $1/C$. A positive answer is
necessary for the route but not sufficient, since the targets must also receive the side information. With the methods
available to us, certifying a single target by the Shannon linear program is limited to about $20$ vertices, while a
triangle-free graph with large $\chi_f$ must be large, so Question~G needs a structural argument in either direction.
